\documentclass[a5paper,10pt]{article}

% https://github.com/InDevRus/filippov-solutions
% Нумерация страниц
\pagenumbering{gobble}

% Настройка полей
\usepackage{geometry}
\geometry{left=1cm}
\geometry{right=1cm}
\geometry{top=1cm}
\geometry{bottom=1.5cm}

% Вставка таблицы
\usepackage{array}
\usepackage{tabularx}

% Инструменты выравнивания формул
\usepackage{amsmath}
\usepackage{mathtools}

% Диакритика в математике
\usepackage{amsfonts}

% Вычисления размеров на ходу
\usepackage{calc}

% Удобные записи производных
\usepackage[italic=true]{derivative}

% Настройки сносок
\usepackage{enotez}

% Длинные знаки векторов
\usepackage{anyfontsize}
\usepackage[b]{esvect}

% Вставка картинок
\usepackage{graphicx}

% Вставка красной строки
\usepackage{indentfirst}
\setlength{\parindent}{1em}

% Условия в командах
\usepackage{xifthen}

% Луа-скрипты
\usepackage{luacode}

% Настройка команд отдельно в математическом и текстовом режимах
\usepackage{mathcommand}

% Для повторения LaTeX кода
\usepackage{multido}

% Конфигурация отступов формул
\delimitershortfall-1sp
\usepackage{mleftright}
\mleftright

% Растягивание объектов
\usepackage{relsize}

% Обтекание текстом
\usepackage{wrapfig}
\usepackage{subcaption}
\usepackage{floatflt}

% Поддержка ввода кириллицы
\usepackage[english, russian]{babel}

% Настройка корневой позиции для компилятора
\usepackage{import}
\providecommand*{\ProjectRootPath}{..}

\import{\ProjectRootPath/preambles/macros/}{theorems}

\import{\ProjectRootPath/preambles/fonts}{font_setup}

% Знак градуса
\usepackage{siunitx}

\import{\ProjectRootPath/preambles/macros/}{operators}
\import{\ProjectRootPath/preambles/macros/}{redefinitions}

\import{\ProjectRootPath/preambles/fonts}{letters_setup}
\import{\ProjectRootPath/preambles/fonts/adjustments}{custom_superscripts}

\import{\ProjectRootPath/preambles/custom}{formatting}
\import{\ProjectRootPath/preambles/custom}{frames}
\import{\ProjectRootPath/preambles/custom}{list_setup}

% TODO: перевести команды на современный стиль объявления
\input{../preambles/plotting}

\begin{document}

$y' = \cos\br{y - x}$. Введём функцию 
$u\br{x} = y\br{x} - x$, $y\br{x} = u\br{x} + x$,
$y'\br{x} = u'\br{x} + 1$.
$u' + 1 = \cos u$.
$u' = \cos u - 1$.
$\dfrac {u'} {\cos u - 1} = 1$, причём решениями являются в том числе функции $u\br{x} = 2\pi n$, где $n$ -- целое число. \vspace{1mm}

$\tint \dfrac {u'\br{x}} {\cos\br{u\br{x}} - 1} dx
= \tint \dfrac {du} {\cos u - 1}
= \begin{bmatrix} u = 2v \end{bmatrix}
= \tint \dfrac {2 dv} {\cos 2v - 1}
= \tint \dfrac {2 dv} {1 - 2\sin^2 v - 1} = \linebreak
= \tint \dfrac {-dv} {\sin^2 v}
= \ctg v + C = \ctg \dfrac {u} {2} + C$. \vspace{1mm}

$\ctg \dfrac {u} {2} = x + C$.
$u\br{x} = 2\arcctg\br{x + C} + 2 \pi n$, где $n$ пробегает ряд целых чисел.

Наконец, общим решением будет
$y\br{x} = x + 2\arcctg\br{x + C} + 2 \pi n$; дополнительные решения равны $y\br{x} = x + 2\pi n$.

\begin{figure}[!h]
	\centering
	\begin{tikzpicture}
		\pgfplotsset{
			Graph/.style = {
				samples = 200,
			},
		}
		\begin{axis}[
			width = 14cm,
			trig format plots = rad,
			ymin = -4.2,
			ymax = 4.2,
			xmin = -5.2,
			xmax = 5.2, 
			xtick = {-10, ..., 10},
			ytick = {-10, ..., 10},
			minor tick num = 3,
			declare function = {
				f(\C, \x) = x + 2 * (pi / 2 - atan(x + \C));
			}]
			
			\pgfplotsinvokeforeach{-0.5, 0, ..., 5} {
				\addplot[Graph, domain = -5.5 : 5.5]
				{f(#1, x)};
			}
			
			\pgfplotsinvokeforeach{-4.5, -4, ..., 0} {
				\addplot[Graph, domain = -5.5 : 5.5]
				{f(#1, x) - 2 * pi};
			}
			
		\end{axis}
	\end{tikzpicture}
\end{figure}

\end{document}
